\section{Finding Top-k Provenance} 
\label{sec:k-provenance}

The two-phase approach described in Section~\ref{sec:1-provenance} returns a single minimal provenance. However, it may not be sufficient for users to verify the correctness of an answer. 
%\agp{move everything from here to TR}
\techreport{For example, consider a question asking for the publication year of a scientific paper in Figure~\ref{fig:top-k}. Provenances 1 and 2 represent two distinct provenances that support the correctness of the answer. In contrast, Provenance 3 happens to return the same answer without supporting its correctness, as it refers a different paper citation in the {\em Reference} section. 

A provenance is considered {\em positive} (e.g., Provenances 1 and 2 in Figure~\ref{fig:top-k}) if it supports the correctness of the answer. Otherwise, it is {\em negative} (e.g., Provenance 3). The distinction between positive and negative provenances can ultimately only be determined by a human than an LLM, as both types of provenance pass the LLM's verification process (i.e., they both reproduce the same answer derived from the full text). } 
To address this issue and to enable users to explore multiple distinct provenances and gain trust in the correctness of the answer, we develop an approach, outlined in Algorithm~\ref{alg:top-k}, to return the $k$ minimal provenances. \techreport{During this process, if users identify any one of the provenances as positive, then the algorithm terminates.} %If all returned provenances are deemed negative by the user, it strongly suggests that the answer is incorrect. 

% Below, we formally define the concepts of sufficient and necessary provenance to formalize the above intuitions. 
% Let $mp.v$ be the label associated with a minimal provenance $mp$, indicating whether $mp$ is positive ($mp.v = \text{P}$) or negative. ($mp.v = \text{N}$)  


% \begin{definition}
%     [Sufficient and Necessary Provenance]~\label{def:sufficient-necessary-provenance} Let $MP = \{mp_1,mp_2,...,mp_n\}$ be the set of minimal provenances within $T$ for the answer $A$ to question $Q$. If $\exists mp_i \in MP$, s.t., $mp_i.v=\text{P}$, then $mp_i$ is a sufficient provenance to prove the answer $A$ is correct. Conversely, if $\forall mp_i \in MP$, $mp_i = \text{N}$, then $MP$ is a necessary provenance set to prove the answer $A$ is incorrect. 
% \end{definition}

% When the answer is correct, a sufficient provenance set containing at least one positive provenance is enough to verify its correctness. Conversely, when the answer is incorrect, the goal is to identify the necessary provenance set, which includes all negative provenances to support this conclusion. In practice, this process may be relaxed by focusing on the top-$k$ provenances to strike a balance between user trust and the algorithm's cost or latency. 

% \begin{example}
%     Consider question answering on the scientific paper in Figure~\ref{fig:example}-2). The top-k provenance approach in \sys (in Section~\ref{sec:k-provenance}) returns one  minimal provenance at a time. If $mp_1$ (Provenance 1) is the first provenance returned, the algorithm terminates, as it is valid and judged by humans to sufficiently demonstrate that the answer is correct. Now, suppose Provenance 3 is the \emph{only} provenance returned by \sys. Since it is invalid, this implies that the answer, "2019," is likely incorrect. 
% \end{example}


%\subsection{Motivation}
%sufficient and necessary provenance 

 




%In contrast to the linear search-based strategies described previously, we present a divide-and-conquer strategy. We begin with the entire set of blocks $B$, where the blocks are unsorted and each block in $B$ preserves the original indexes from the text $T$. In each iteration, we divide the current list of blocks $T_{l,r}$ into a left half $T_{l,mid}$ and a right half $T_{mid+1,r}$, where $mid = \lceil \frac{l + r}{2} \rceil$. We maintain sibling relationships between $T_{l,mid}$ and $T_{mid+1,r}$, as well as parent relationships for $T_{l,mid}$, $T_{mid+1,r}$, and $T_{l,r}$ accordingly (Line 9–11). The strategy recursively searches the half that returns an answer equivalent to $A$ (Line 13-14). The strategy terminates when only one block remains (i.e., $l$ equals $r$ for $T_{l,r}$) (Line 3-5) or, when both the current blocks $T_{l,r}$ and its sibling blocks $SIB(T_{l,r})$ evaluate as false, where in such a case the parent blocks $PAR(T_{l,r})$ is returned as provenance (Line 6-9).   

% \begin{algorithm}[tb]
% \small
% \caption{\code{top-k-provenance}}
% \label{alg:top-k}
% \KwIn{$Q,\,T,\,L,\,A,F,I$}  
% %$P \leftarrow \emptyset$ \\ 
% $T = \langle b_1,\dots ,b_m\rangle$ \\ 
% $H \leftarrow \text{empty max heap}$\\
% $H.append((T_{1,m}, 1)) $ \\  
% \While{$H \neq \emptyset$}{
% $(T_{l,r}, score_{l,r}) \leftarrow H.pop()$ \\ 
% \If{$l < r$}{
% $mid = \lceil \frac{l+r}{2} \rceil$ \\ 
%     $score_{l,mid} \leftarrow F(L(T_{l,mid},Q),A)$ \\ 
%     $score_{mid+1,r} \leftarrow F(L(T_{mid+1,r},Q),A)$ \\ 
%      }
% $c_1 \leftarrow (l == r) \land (score_{l,r} \geq 0.5)$ \\ 
% $c_2 \leftarrow (l < r) \land (score_{l,r} \geq 0.5) \land $  $(score_{l,mid} < 0.5) \land $  $(score_{mid+1,r} < 0.5)$ \\ 
% \If{$c_2 == False$}{
% $H.append((T_{l,mid}, score_{l,mid}))$ \\ 
% $H.append((T_{mid+1,r}, score_{mid+1,r}))$ \\
% }
%     \If{$c_1 == True$ or $c_2 == True$}{
%     \If{$c_1 == True$}{
%     $P_{cur} = \code{refine}(T_{l,r})$ \\ }
%     \Else{$P_{cur} = \code{refine(prune(}(T_{l,r})$)) \\}
%     \techreport{\If{\code{user\_verify}($P_{cur}$) == Positive}{
%         {\bf Return} $P_{cur}$, True 
%     }}
%     %$P \leftarrow P \cup P_{cur}$ \\ 
%     \If{$|P| \geq k$}{{\bf Return } $P_{cur}$, False}
%     }
% }
% \end{algorithm}

\begin{algorithm}[tb]
\scriptsize
\caption{\code{top-k-provenance}}
\label{alg:top-k}
\KwIn{$Q,\,T,\,L,\,A,F,I,\rtwo{k}$}  
%$P \leftarrow \emptyset$ \\ 
$T = \langle b_1,\dots ,b_m\rangle$ \\ 
$H \leftarrow \text{empty max heap}$, $P \leftarrow \text{empty list}$ \\ 
$H.append((T_{1,m}, 1)) $ \\  
\While{$H \neq \emptyset$}{
$(T_{l,r}, score_{l,r}) \leftarrow H.pop()$ \\ 
\If{$l < r$}{
$mid = \lceil \frac{l+r}{2} \rceil$ \\ 
    $score_{l,mid} \leftarrow F(L(T_{l,mid},Q),A)$ \\ 
    $score_{mid+1,r} \leftarrow F(L(T_{mid+1,r},Q),A)$ \\ 
     }
$\rtwo{c_1 \leftarrow (l == r) \land (score_{l,r} \geq 0.5)}$ \\ 
\rtwo{$c_2 \leftarrow (l < r) \land (score_{l,r} \geq 0.5)  \land (score_{l,mid} < 0.5) \land   (score_{mid+1,r} < 0.5)$} \\
\If{\rtwo{$(c_1 \lor c_2) == True$}}{
    \If{\rtwo{$c_1 == True$}}{
    $P_{cur} = \code{refine}(T_{l,r})$ \\ }
    \Else{$P_{cur} = \code{refine(prune(}(T_{l,r})$)) \\}
    \techreport{\If{\code{user\_verify}($P_{cur}$) == Positive}{
        {\bf Return} $P_{cur}$, True 
    }}
    $P \leftarrow P \cup P_{cur}$ \\ 
}
\If{$c_2 == False$}{
$H.append((T_{l,mid}, score_{l,mid}))$ \\ 
$H.append((T_{mid+1,r}, score_{mid+1,r}))$ \\
}
\If{$|P| \geq k$}{{\bf Return } $P$}
    }
\end{algorithm}



\topic{Algorithm} 
\rtwo{We illustrate Algorithm~\ref{alg:top-k} for inferring $k$ provenances using a running example shown in Figure~\ref{fig:top-k}. We let the text $T$ be a list of $m$ equal-sized blocks, and we aim to find three minimal provenances ($k=3$).}

\rtwo{Initially, we start from $T$ (or equivalently, $T_{1,m}$, denoting the sequence of blocks from $B_1$ to $B_m$), represented by the root node 1. In the second iteration, we split $T_{1,m}$ into left and right halves, corresponding to its child node 2 $T_{1,\lceil m/2 \rceil}$ and node 3 $T_{\lceil m/2 \rceil + 1,m}$, respectively.} For each node, we compute a score as the likelihood that the corresponding block can reproduce $A$. Such a score is 1 if the answer $A'$ produced by the block is lexically identical to $A$; otherwise, we leverage LLMs to generate the score \papertext{(see prompt in \vldb{\cite{blip}}). }\techreport{using the prompt \code{Top-k-Eval-Prompt} below.}   
\rtwo{For example, the scores of nodes 1 and 2 are 1 and 0.9.} 

\techreport{\noindent{\underline{\code{Top-k-Eval-Prompt}}}: \textcolor{blue}{[Question]}, \textcolor{blue}{[Answer 1]}, and \textcolor{blue}{[Answer 2]} are placeholders for the question $Q$, the answer $L(T_{l,\text{mid}}, Q)$ evaluated on $T_{l,\text{mid}}$, and $A$, respectively, in Line 8-9 in Algorithm~\ref{alg:top-k}.  Note that when \textcolor{blue}{value} is \textit{False}, we use $1 -$~\textcolor{blue}{score} as its likelihood. 


\mypython{
\small 
    Top-k-Eval-Prompt: Given the following question, \textcolor{blue}{[Question]}, and two answers, \textcolor{blue}{[Answer 1]} and \textcolor{blue}{[Answer 2]}. Determine whether the two answers are equivalent in meaning. Return the result as a JSON object with the following format: \textcolor{blue}{value}: true if the answers are equivalent, false otherwise. \textcolor{blue}{score}: a real number between 0 and 1 representing the likelihood that your judgment is correct. }}


\techreport{
\begin{figure}[]
    \centering
    \includegraphics[width=0.9\linewidth]{figures/paper_example.pdf}
    %\vspace{-2em}
    \caption{\small Top-k Provenance. }  
    %\vspace{-1em}
    \label{fig:top-k}
\end{figure}
}

We create a max heap $H$ that stores pairs consisting of a text block and its corresponding score. In each iteration, we expand the node with the highest likelihood score from the max heap and remove the node that is already expanded.  
\rtwo{For example, node 1 is removed from $H$ in the second iteration, and node 2 is expanded in the next iteration since it has the highest score.} We then insert its two children (node 4 and 5)—obtained by splitting the current text sequence into left and right halves—into the max heap and rank them based on their likelihood scores.

\rtwo{The above node splitting continues recursively until either of the following two conditions is met.} \rtwo{First, when we reach a node that contains only one block (e.g., node 8) and it passes verification (i.e., reproduces the answer), we return it by invoking \code{refine} to obtain the corresponding minimal provenance.} In this example, node 8 is the first provenance returned by the \code{top-k-provenance} algorithm\papertext{.}\techreport{, assuming it is distinct from the provenance returned by the two-phase approach.}  
\rtwo{Second, when we reach an internal node (e.g., node 6) whose score is no less than 0.5, but both of its child nodes have scores lower than 0.5 (e.g., nodes 10 and 11), node 6 is returned as the third provenance. This provenance is then further refined using the two-phase approach to obtain a minimal provenance.}



% Consider the text $T$ as a list of $m$ equal-sized blocks (Line 2). We create a max heap $H$ that stores pairs consisting of a text block and its corresponding score, denoting the likelihood that this block can reproduce the answer, as will be explained shortly. Let $T_{i,j} = \langle B_i, \dots, B_j \rangle$ denote the sequence of blocks from $B_i$ to $B_j$. We initialize $H$ with the pair $(T_{1,m}, 1)$ (Line 2-3). The max heap $H$ sorts the pairs based on their scores in descending order. In each iteration, we extract the block sequence $T_{l,r}$ with the highest score from $H$ and split it into left and right halves, i.e., $T_{l,\text{mid}}$ and $T_{\text{mid}+1,r}$, where $\text{mid} = \lceil \frac{l + r}{2} \rceil$ (Line 5-7). 

% For each half of the block sequence, such as  $T_{l,\text{mid}}$, we generate an answer $L(T_{l,\text{mid}}, Q)$ to the question $Q$ using the same LLM $L$, and apply a scoring function $F$ to compute a likelihood score $score_{l,\text{mid}}$ indicating how likely the generated answer $L(T_{l,\text{mid}}, Q)$ is equivalent to $A$. If $L(T_{l,\text{mid}}, Q)$ is lexically identical to $A$, the likelihood score is set to 1. Otherwise, we leverage LLMs to generate the score using the \code{Top-k-Eval-Prompt} in \sigmod{Appendix~\ref{subsec:appendix-prompt}}\vldb{\cite{blip}} (Line 8-9). 


% The algorithm returns a provenance when either of the below two conditions is met. First, when a sequence consisting of a single block ($l == r$) produces an answer equivalent to $A$ (indicated by $score_{l,r} \geq 0.5$ in Line 10), we refine this block using the \code{refine} from Algorithm~\ref{alg:greedy} to obtain a minimal provenance (Lines 16–17). Second, when a sequence $T_{l,r}$ ($l<r$) can reproduce the answer\papertext{,} \techreport{(with a score of at least 0.5),} but neither of its halves, i.e., $T_{l,\text{mid}}$ nor $T_{\text{mid}+1,r}$, can (i.e., $c_2$ in Line 11), we stop exploring its subsequences and instead return $T_{l,r}$ as a candidate. We then invoke the two-phase approach to produce a minimal provenance from it (Lines 18–19). A \texttt{True} value for $c_2$ suggests that the provenance likely includes information from both $T_{l,\text{mid}}$ and $T_{\text{mid}+1,r}$, and that using only one of them is insufficient to form a provenance. 
% \techreport{If the user verifies the returned minimal provenance as positive, the algorithm terminates. Otherwise, the}\papertext{The} search continues for additional minimal provenances until either $k$ provenances are returned or the max heap becomes empty, indicating that all possible minimal provenances have been explored \techreport{(Line 20-23).}\papertext{(Line 20-21).}  

% Empirically, we find that the top-k algorithm
% produces provenances at a small additional cost
% relative to one that produces a single minimal
% provenance. 


% \begin{example}
    % %example
% Each node in the tree represents a sequence of text blocks. The root node corresponds to $T_{1,m}$, and its child nodes 2 and 3 represent $T_{1,\lceil m/2 \rceil}$ and $T_{\lceil m/2 \rceil + 1,m}$, respectively. 
% Assume we aim to identify the top-3 provenances ($k=3$). In each iteration, we expand the node with the highest likelihood score from the max heap corresponding to the leftmost value in each row of the table in Figure~\ref{fig:top-k}. We then insert its two children—obtained by splitting the current text sequence into left and right halves—into the max heap, and rank them based on their likelihood scores. 
% When we reach a leaf node that contains only one block (e.g., node 8) and it passes the verification (i.e., reproduces the answer), we return it by invoking \code{refine} to obtain the corresponding minimal provenance. In this example, node 8 is the first provenance returned by the \code{top-k-provenance} algorithm\papertext{.}\techreport{, assuming it is distinct from the provenance returned by the two-phase approach.} When we reach an internal node (e.g., node 6) whose score is no less than 0.5, but both of its child nodes have scores lower than 0.5 (e.g., nodes 10 and 11), node 6 is returned as a provenance. This provenance is then further refined using the two-phase approach to obtain a minimal provenance. 
% %end of example 
% \end{example}



\techreport{\topic{Discussion} %case 1: which scenarios top-k-provenance can return all provenances? 
% case 2: if not, then what is the argument? one way can try to explore is: if there exist positive provenances, our algorithm can discover at least one. Then it would be sufficient to remove the need of finding all minimal provenances. 
The \code{top-k-provenance} algorithm does not guarantee returning all distinct minimal provenances, but only a subset of them in $T$. We argue that this is often sufficient in many scenarios. If the question over $T$ is correctly answered by LLM $L$, then finding a single positive provenance suffices for user trust, and the algorithm can terminate without exploring all minimal provenances. On the other hand, if the answer returned by $L$ on $T$ is incorrect, ideally  all minimal provenances that are negative are expected to trust that answer is wrong. }

\begin{figure}[tb]
    \centering
    \vspace{-1mm}
    \includegraphics[width=1\linewidth]{figures/topk1.pdf}
    \vspace{-8mm}
    \caption{\small Top-$k$ Provenance.} 
    \vspace{-1mm}
    \label{fig:top-k} 
\end{figure}

%In practice, if none of the provenances ranked by their likelihood of being valid from \code{top-k-provenance} is positive, it implies a higher chance that the answer is incorrect. In such cases, users may not wish to explore all minimal provenances due to the non-trivial cost and latency involved. 
 




% \subsection{Analysis}

% \subsubsection{Cost Analysis}
% \leavevmode
% %KV-cache 


% \subsubsection{Result Completeness}
% \leavevmode


 